% --- Archon physics formalization source begin ---
% archon:physics
% archon:covers PhyXMiniProblems/problem_phyx_mini_0402.lean
% archon:source-report reports/phyx_mini/problem_phyx_mini_0402.source.json

\chapter{Physics problem phyx\_mini\_0402}
\label{ch:PhyXMiniProblems_problem_phyx_mini_0402}

\paragraph{Problem source.}
\#\# Physical scenario

A gas is compressed in two stages. First, its volume decreases while pressure rises. Then, its volume further decreases as pressure declines, resulting in a smaller final volume.

\#\# Question

How much work is done on the gas in the process shown in figure?

\#\# Answer choices

- A: A: 0J

- B: B: -60J

- C: C: 60J

- D: D: 30

\#\# Figure caption (auxiliary; use the image as primary evidence)

The image is a graph with pressure (p) in kilopascals (kPa) on the vertical axis and volume (V) in cubic centimeters (cm³) on the horizontal axis. The graph shows a plot with two notable points labeled "f" and "i." 

- A point labeled "f" is located at approximately 90 cm³ and 200 kPa.
- A point labeled "i" is located at around 285 cm³ and 200 kPa.
- The plot consists of two line segments: 
  - An upward-sloping line from point "f" to a peak.
  - A downward-sloping line from the peak to point "i." 
- Each segment is marked with arrows indicating the direction of movement along the plot. 
- The peak between the two points is around 190 cm³ and 400 kPa. 

The graph appears to depict a thermodynamic cycle or process involving changes in pressure and volume.

\#\# Dataset metadata

Laws of Thermodynamics; Physical Model Grounding Reasoning, Multi-Formula Reasoning

\paragraph{Recorded answer/context.}
C: C: 60J

\paragraph{Figure/image path.}
/root/proposal\_for\_physic/science-mango/phyx\_mini\_run/phyx\_data/test\_image/402.png

\paragraph{Formalization target.}
create a compiling Lean file with sorry bodies at `PhyXMiniProblems/problem\_phyx\_mini\_0402.lean`. The Lean declarations must preserve the physical quantities, dimensions or dimensional roles, figure labels, governing-law hypotheses, and final relation expressed by this problem.
Use Mathlib/Physlib names found through LeanExplore where available. If a physics API is missing, introduce faithful local abstractions rather than scalar placeholder aliases.

\begin{theorem}[Physics formalization target]
\label{thm:physics:phyx_mini_0402:target}
\lean{PhyXMiniProblems.ProblemPhyXMini0402.work_done_on_gas_is_sixty_joules}
\uses{def:physics:phyx-mini-0402:phyxminiproblems-problemphyxmini0402-energyinjoules, def:physics:phyx-mini-0402:phyxminiproblems-problemphyxmini0402-twostagegascompression, def:physics:phyx-mini-0402:phyxminiproblems-problemphyxmini0402-matchestwostagecompressionscenario, def:physics:phyx-mini-0402:phyxminiproblems-problemphyxmini0402-matchesprimarypressurevolumefigure, def:physics:phyx-mini-0402:phyxminiproblems-problemphyxmini0402-obeyspiecewiselinearpressurevolumeworklaw}
For the two straight compression legs shown in the primary pressure--volume
figure, the total work done on the gas is \(60\) joules.
\end{theorem}
\begin{proof}
The plotted states are
\[
 (300\ {\rm cm^3},200\ {\rm kPa}),\quad
 (200\ {\rm cm^3},400\ {\rm kPa}),\quad
 (100\ {\rm cm^3},200\ {\rm kPa}).
\]
On a straight segment the work on the gas is the average endpoint pressure
times the decrease in volume.  Each leg therefore contributes
\[
 \frac{200+400}{2}\,(300-200)\,10^{-3}=30\ {\rm J}
\]
and
\[
 \frac{400+200}{2}\,(200-100)\,10^{-3}=30\ {\rm J},
\]
respectively.  The total-work law adds the two contributions, giving
\(60\) joules.
\end{proof}

% --- Archon declaration topology begin ---
\section{Declaration topology}
\begin{definition}[Gas Volume]
  \label{def:physics:phyx-mini-0402:phyxminiproblems-problemphyxmini0402-gasvolume}
  \lean{PhyXMiniProblems.ProblemPhyXMini0402.GasVolume}
  A nonnegative physical gas volume, carrying dimension L³
\end{definition}
\begin{proof}
  This is the declared model definition of Gas Volume.
\end{proof}
\begin{definition}[volume In Cubic Centimeters]
  \label{def:physics:phyx-mini-0402:phyxminiproblems-problemphyxmini0402-volumeincubiccentimeters}
  \lean{PhyXMiniProblems.ProblemPhyXMini0402.volumeInCubicCentimeters}
  \uses{def:physics:phyx-mini-0402:phyxminiproblems-problemphyxmini0402-gasvolume}
  Read a physical gas volume in cubic centimetres
\end{definition}
\begin{proof}
  This is the declared model definition of volume In Cubic Centimeters.
\end{proof}
\begin{definition}[pressure In Pascals]
  \label{def:physics:phyx-mini-0402:phyxminiproblems-problemphyxmini0402-pressureinpascals}
  \lean{PhyXMiniProblems.ProblemPhyXMini0402.pressureInPascals}
  Read a physical pressure in pascals
\end{definition}
\begin{proof}
  This is the declared model definition of pressure In Pascals.
\end{proof}
\begin{definition}[pressure In Kilopascals]
  \label{def:physics:phyx-mini-0402:phyxminiproblems-problemphyxmini0402-pressureinkilopascals}
  \lean{PhyXMiniProblems.ProblemPhyXMini0402.pressureInKilopascals}
  \uses{def:physics:phyx-mini-0402:phyxminiproblems-problemphyxmini0402-pressureinpascals}
  Read a physical pressure in kilopascals, the figure's vertical unit
\end{definition}
\begin{proof}
  This is the declared model definition of pressure In Kilopascals.
\end{proof}
\begin{definition}[energy In Joules]
  \label{def:physics:phyx-mini-0402:phyxminiproblems-problemphyxmini0402-energyinjoules}
  \lean{PhyXMiniProblems.ProblemPhyXMini0402.energyInJoules}
  Read a signed physical work or energy quantity in joules
\end{definition}
\begin{proof}
  This is the declared model definition of energy In Joules.
\end{proof}
\begin{definition}[joules Per Kilopascal Cubic Centimeter]
  \label{def:physics:phyx-mini-0402:phyxminiproblems-problemphyxmini0402-joulesperkilopascalcubiccentimeter}
  \lean{PhyXMiniProblems.ProblemPhyXMini0402.joulesPerKilopascalCubicCentimeter}
  One kPa · cm³ is 10⁻³ J
\end{definition}
\begin{proof}
  This is the declared model definition of joules Per Kilopascal Cubic Centimeter.
\end{proof}
\begin{definition}[straight Segment Work On Gas In Joules]
  \label{def:physics:phyx-mini-0402:phyxminiproblems-problemphyxmini0402-straightsegmentworkongasinjoules}
  \lean{PhyXMiniProblems.ProblemPhyXMini0402.straightSegmentWorkOnGasInJoules}
  \uses{def:physics:phyx-mini-0402:phyxminiproblems-problemphyxmini0402-joulesperkilopascalcubiccentimeter}
  Work done on the gas along one straight segment of a p--V diagram. The average pressure on a straight segment is the arithmetic mean of its endpoint pressures. The factor initialVolume - finalVolume implements the work-on-the-gas sign convention, so a compression has positive work
\end{definition}
\begin{proof}
  This is the declared model definition of straight Segment Work On Gas In Joules.
\end{proof}
\begin{definition}[Diagram Point]
  \label{def:physics:phyx-mini-0402:phyxminiproblems-problemphyxmini0402-diagrampoint}
  \lean{PhyXMiniProblems.ProblemPhyXMini0402.DiagramPoint}
  The two labelled endpoints and the unlabelled peak in the image
\end{definition}
\begin{proof}
  This is the declared model definition of Diagram Point.
\end{proof}
\begin{definition}[Process Segment]
  \label{def:physics:phyx-mini-0402:phyxminiproblems-problemphyxmini0402-processsegment}
  \lean{PhyXMiniProblems.ProblemPhyXMini0402.ProcessSegment}
  The two directed straight segments shown by arrows in the image
\end{definition}
\begin{proof}
  This is the declared model definition of Process Segment.
\end{proof}
\begin{definition}[initial Point]
  \label{def:physics:phyx-mini-0402:phyxminiproblems-problemphyxmini0402-processsegment-initialpoint}
  \lean{PhyXMiniProblems.ProblemPhyXMini0402.ProcessSegment.initialPoint}
  \uses{def:physics:phyx-mini-0402:phyxminiproblems-problemphyxmini0402-diagrampoint, def:physics:phyx-mini-0402:phyxminiproblems-problemphyxmini0402-processsegment}
  Initial diagram point of each directed segment
\end{definition}
\begin{proof}
  This is the declared model definition of initial Point.
\end{proof}
\begin{definition}[final Point]
  \label{def:physics:phyx-mini-0402:phyxminiproblems-problemphyxmini0402-processsegment-finalpoint}
  \lean{PhyXMiniProblems.ProblemPhyXMini0402.ProcessSegment.finalPoint}
  \uses{def:physics:phyx-mini-0402:phyxminiproblems-problemphyxmini0402-diagrampoint, def:physics:phyx-mini-0402:phyxminiproblems-problemphyxmini0402-processsegment}
  Final diagram point of each directed segment
\end{definition}
\begin{proof}
  This is the declared model definition of final Point.
\end{proof}
\begin{definition}[Axis Quantity]
  \label{def:physics:phyx-mini-0402:phyxminiproblems-problemphyxmini0402-axisquantity}
  \lean{PhyXMiniProblems.ProblemPhyXMini0402.AxisQuantity}
  Physical quantity assigned to an axis in the supplied figure
\end{definition}
\begin{proof}
  This is the declared model definition of Axis Quantity.
\end{proof}
\begin{definition}[Volume Axis Unit]
  \label{def:physics:phyx-mini-0402:phyxminiproblems-problemphyxmini0402-volumeaxisunit}
  \lean{PhyXMiniProblems.ProblemPhyXMini0402.VolumeAxisUnit}
  Unit printed on the horizontal axis
\end{definition}
\begin{proof}
  This is the declared model definition of Volume Axis Unit.
\end{proof}
\begin{definition}[Pressure Axis Unit]
  \label{def:physics:phyx-mini-0402:phyxminiproblems-problemphyxmini0402-pressureaxisunit}
  \lean{PhyXMiniProblems.ProblemPhyXMini0402.PressureAxisUnit}
  Unit printed on the vertical axis
\end{definition}
\begin{proof}
  This is the declared model definition of Pressure Axis Unit.
\end{proof}
\begin{definition}[Segment Shape]
  \label{def:physics:phyx-mini-0402:phyxminiproblems-problemphyxmini0402-segmentshape}
  \lean{PhyXMiniProblems.ProblemPhyXMini0402.SegmentShape}
  Geometric shape of each process leg in the p--V plane
\end{definition}
\begin{proof}
  This is the declared model definition of Segment Shape.
\end{proof}
\begin{definition}[Process Kind]
  \label{def:physics:phyx-mini-0402:phyxminiproblems-problemphyxmini0402-processkind}
  \lean{PhyXMiniProblems.ProblemPhyXMini0402.ProcessKind}
  Thermodynamic role of each leg
\end{definition}
\begin{proof}
  This is the declared model definition of Process Kind.
\end{proof}
\begin{definition}[Pressure Trend]
  \label{def:physics:phyx-mini-0402:phyxminiproblems-problemphyxmini0402-pressuretrend}
  \lean{PhyXMiniProblems.ProblemPhyXMini0402.PressureTrend}
  Qualitative pressure change along a directed leg
\end{definition}
\begin{proof}
  This is the declared model definition of Pressure Trend.
\end{proof}
\begin{definition}[Thermodynamic State]
  \label{def:physics:phyx-mini-0402:phyxminiproblems-problemphyxmini0402-thermodynamicstate}
  \lean{PhyXMiniProblems.ProblemPhyXMini0402.ThermodynamicState}
  \uses{def:physics:phyx-mini-0402:phyxminiproblems-problemphyxmini0402-gasvolume}
  Physical pressure and volume at one point of the process
\end{definition}
\begin{proof}
  This is the declared model definition of Thermodynamic State.
\end{proof}
\begin{definition}[Pressure Volume Figure]
  \label{def:physics:phyx-mini-0402:phyxminiproblems-problemphyxmini0402-pressurevolumefigure}
  \lean{PhyXMiniProblems.ProblemPhyXMini0402.PressureVolumeFigure}
  \uses{def:physics:phyx-mini-0402:phyxminiproblems-problemphyxmini0402-diagrampoint, def:physics:phyx-mini-0402:phyxminiproblems-problemphyxmini0402-processsegment, def:physics:phyx-mini-0402:phyxminiproblems-problemphyxmini0402-axisquantity, def:physics:phyx-mini-0402:phyxminiproblems-problemphyxmini0402-volumeaxisunit, def:physics:phyx-mini-0402:phyxminiproblems-problemphyxmini0402-pressureaxisunit, def:physics:phyx-mini-0402:phyxminiproblems-problemphyxmini0402-segmentshape}
  The axes, plotted coordinates, labels, segments, and arrows visible in the primary image. Coordinates are scalar readouts in the units stored alongside the axes; they are not replacements for physical pressure or volume
\end{definition}
\begin{proof}
  This is the declared model definition of Pressure Volume Figure.
\end{proof}
\begin{definition}[Two Stage Gas Compression]
  \label{def:physics:phyx-mini-0402:phyxminiproblems-problemphyxmini0402-twostagegascompression}
  \lean{PhyXMiniProblems.ProblemPhyXMini0402.TwoStageGasCompression}
  \uses{def:physics:phyx-mini-0402:phyxminiproblems-problemphyxmini0402-diagrampoint, def:physics:phyx-mini-0402:phyxminiproblems-problemphyxmini0402-processsegment, def:physics:phyx-mini-0402:phyxminiproblems-problemphyxmini0402-processkind, def:physics:phyx-mini-0402:phyxminiproblems-problemphyxmini0402-pressuretrend, def:physics:phyx-mini-0402:phyxminiproblems-problemphyxmini0402-thermodynamicstate, def:physics:phyx-mini-0402:phyxminiproblems-problemphyxmini0402-pressurevolumefigure}
  A gas sample and its two-stage compression. The gas carrier is abstract, while all plotted state variables and work values retain physical dimensions. The segment works and total work are independent fields; no requested numerical answer is assigned here
\end{definition}
\begin{proof}
  This is the declared model definition of Two Stage Gas Compression.
\end{proof}
\begin{definition}[Matches Two Stage Compression Scenario]
  \label{def:physics:phyx-mini-0402:phyxminiproblems-problemphyxmini0402-matchestwostagecompressionscenario}
  \lean{PhyXMiniProblems.ProblemPhyXMini0402.MatchesTwoStageCompressionScenario}
  \uses{def:physics:phyx-mini-0402:phyxminiproblems-problemphyxmini0402-twostagegascompression}
  The prose classification of the two legs. It contains no numerical work claim
\end{definition}
\begin{proof}
  This is the declared model definition of Matches Two Stage Compression Scenario.
\end{proof}
\begin{definition}[Matches Primary Pressure Volume Figure]
  \label{def:physics:phyx-mini-0402:phyxminiproblems-problemphyxmini0402-matchesprimarypressurevolumefigure}
  \lean{PhyXMiniProblems.ProblemPhyXMini0402.MatchesPrimaryPressureVolumeFigure}
  \uses{def:physics:phyx-mini-0402:phyxminiproblems-problemphyxmini0402-volumeincubiccentimeters, def:physics:phyx-mini-0402:phyxminiproblems-problemphyxmini0402-pressureinkilopascals, def:physics:phyx-mini-0402:phyxminiproblems-problemphyxmini0402-twostagegascompression}
  Exact transcription of the primary pressure-volume raster, including the axis roles and units, the three plotted coordinates, endpoint labels, straight segments, and arrow directions. The physical state readouts are explicitly calibrated to these coordinates
\end{definition}
\begin{proof}
  This is the declared model definition of Matches Primary Pressure Volume Figure.
\end{proof}
\begin{definition}[Obeys Piecewise Linear Pressure Volume Work Law]
  \label{def:physics:phyx-mini-0402:phyxminiproblems-problemphyxmini0402-obeyspiecewiselinearpressurevolumeworklaw}
  \lean{PhyXMiniProblems.ProblemPhyXMini0402.ObeysPiecewiseLinearPressureVolumeWorkLaw}
  \uses{def:physics:phyx-mini-0402:phyxminiproblems-problemphyxmini0402-volumeincubiccentimeters, def:physics:phyx-mini-0402:phyxminiproblems-problemphyxmini0402-pressureinkilopascals, def:physics:phyx-mini-0402:phyxminiproblems-problemphyxmini0402-energyinjoules, def:physics:phyx-mini-0402:phyxminiproblems-problemphyxmini0402-straightsegmentworkongasinjoules, def:physics:phyx-mini-0402:phyxminiproblems-problemphyxmini0402-twostagegascompression}
  For each straight leg, work on the gas is average pressure times the decrease in volume, with the explicit kPa · cm³ to joule conversion. The total work is the sum of the two independent leg works. These laws are generic in the state data and do not contain the requested 60 J conclusion
\end{definition}
\begin{proof}
  This is the declared model definition of Obeys Piecewise Linear Pressure Volume Work Law.
\end{proof}
% --- Archon declaration topology end ---
% --- Archon physics formalization source end ---
